% @@CORSO@@: appunti
% Compilare con LuaLaTeX: latexmk -lualatex main.tex
\documentclass[11pt,a4paper,openany]{book}

\usepackage[a4paper,margin=2.6cm,headheight=14pt]{geometry}

% --- Lingue ---------------------------------------------------------------------
% Aggiungi le lingue del materiale prima di main=, per esempio:
%   \usepackage[ngerman,latin,french,spanish,english,main=italian]{babel}
% Il greco, anche politonico, si scrive direttamente con Libertinus: non serve dichiararlo.
\usepackage[english,main=italian]{babel}
\babelfont{rm}{Libertinus Serif}
\babelfont{sf}{Libertinus Sans}
\babelfont{tt}{Libertinus Mono}
% Per l'ebraico aggiungi bidi=basic alle opzioni di babel e togli il commento a queste righe:
% \babelprovide[import]{hebrew}
% \babelfont[hebrew]{rm}[Renderer=HarfBuzz]{Libertinus Serif}

% --- Pacchetti --------------------------------------------------------------------
\usepackage{amsmath}
\usepackage{unicode-math}
\setmathfont{Libertinus Math}
\usepackage{microtype}
\usepackage{graphicx}
\graphicspath{{images/}}
\usepackage{float}
\usepackage{etoolbox}
\usepackage[font=small,labelfont=bf]{caption}
\usepackage{booktabs,tabularx,array,longtable,fvextra}
\usepackage{ragged2e}
\newcolumntype{Y}{>{\RaggedRight\arraybackslash}X}
\usepackage{enumitem}
\setlist{itemsep=2pt,topsep=4pt}
\usepackage{paracol}
\usepackage{xcolor}
\usepackage[most]{tcolorbox}
\usepackage{listings}
\usepackage{tikz}
\usetikzlibrary{arrows.meta,positioning,shapes.geometric,automata,babel}
\usepackage{imakeidx}
\makeindex[name=indice,title=Indice analitico,intoc,columns=2]
\usepackage{fancyhdr}
\usepackage[hidelinks,unicode]{hyperref}

\definecolor{accento}{HTML}{1F4E79}
\definecolor{corr}{HTML}{A4262C}
\definecolor{libro}{HTML}{1F6F5C}
\definecolor{fonte}{HTML}{7A5C00}
\definecolor{sintesi}{HTML}{4A4A6A}

\pagestyle{fancy}
\fancyhf{}
\fancyhead[LE,RO]{\small\thepage}
\fancyhead[RE]{\small\ultimocapitolo}
\fancyhead[LO]{\small\ultimasezione}
\renewcommand{\headrulewidth}{0.3pt}
% Le intestazioni leggono queste macro e non \leftmark e \rightmark, che con paracol
% restano vuoti o arretrati di un capitolo.
\makeatletter
\newcommand{\ultimocapitolo}{}
\newcommand{\ultimasezione}{}
\renewcommand{\chaptermark}[1]{%
  \if@mainmatter\xdef\ultimocapitolo{\@chapapp\ \thechapter.\noexpand\ \unexpanded{#1}}%
  \else\gdef\ultimocapitolo{#1}\fi
  \gdef\ultimasezione{}\markboth{\ultimocapitolo}{\ultimasezione}}
\renewcommand{\sectionmark}[1]{%
  \xdef\ultimasezione{\thesection.\noexpand\ \unexpanded{#1}}\markright{\ultimasezione}}
\makeatother

% Codice letterale: conserva spazi e Unicode; spezza anche percorsi senza spazi.
% lstlisting usa le opzioni di fvextra (per esempio numbers=left), non language=C.
\DefineVerbatimEnvironment{lstlisting}{Verbatim}{breaklines=true,breakanywhere=true,
  breakautoindent=false,breakindent=1em,breaksymbolleft={},breaksymbolright={},
  fontsize=\normalsize,frame=leftline,framesep=5pt}
\renewcommand{\theFancyVerbLine}{\normalfont\normalsize\arabic{FancyVerbLine}}


% --- Riquadri -----------------------------------------------------------------------
\tcbset{appbox/.style={enhanced,breakable,boxrule=0pt,leftrule=2.5pt,sharp corners,
  left=8pt,right=8pt,top=5pt,bottom=5pt,before skip=9pt,after skip=9pt,
  fonttitle=\bfseries\small,coltitle=black,attach title to upper={\par\smallskip},
  fontupper=\small}}
% Errore o incongruenza nel materiale: il testo originale resta visibile.
\newtcolorbox{correzione}[1][]{appbox,colframe=corr,colback=corr!6,title={\textcolor{corr}{Correzione}},#1}
% Dubbio non risolvibile con le fonti del corso.
\newtcolorbox{dubbio}[1][]{appbox,colframe=fonte,colback=fonte!7,title={\textcolor{fonte}{Da verificare}},#1}
% Spiegazione tratta dal libro, con la pagina.
\newtcolorbox{dallibro}[1][]{appbox,colframe=libro,colback=libro!7,title={\textcolor{libro}{Dal libro}},#1}
% Sintesi concettuale a fine capitolo.
\newtcolorbox{insintesi}[1][]{appbox,colframe=sintesi,colback=sintesi!7,title={\textcolor{sintesi}{In sintesi}},#1}

% Sezione senza slide di riferimento.
\newcommand{\soloappunti}{\par\noindent\textcolor{fonte}{\small\sffamily\bfseries Fonte: solo appunti in classe (nessuna slide)}\par\smallskip}
% Aggiunta di conoscenza generale, non dal materiale del corso.
\newcommand{\integrazione}[1]{\textcolor{sintesi}{\small[\emph{Integrazione, non dal materiale del corso:} #1]}}

% --- Citazioni e testi a fronte -----------------------------------------------------
% La lingua facoltativa deve essere dichiarata in babel.
\newenvironment{cit}[1][italian]{\begin{quote}\small\selectlanguage{#1}}{\end{quote}}
\newcommand{\fontecit}[1]{\par\hfill{\footnotesize\selectlanguage{italian}#1}}
\newcommand{\afronte}[3][italian]{%
  \begin{paracol}{2}\small\selectlanguage{#1}\sloppy #2\switchcolumn\selectlanguage{italian}\sloppy #3\end{paracol}\medskip}

% --- Scritture e parole straniere ----------------------------------------------------
\newcommand{\gr}[1]{#1}
\newcommand{\en}[1]{\foreignlanguage{english}{\emph{#1}}}
% Togli il commento dopo aver dichiarato la lingua in babel:
% \newcommand{\de}[1]{\foreignlanguage{ngerman}{\emph{#1}}}
% \newcommand{\la}[1]{\foreignlanguage{latin}{\emph{#1}}}
% \newcommand{\fr}[1]{\foreignlanguage{french}{\emph{#1}}}
% \newcommand{\es}[1]{\foreignlanguage{spanish}{\emph{#1}}}
% \newcommand{\he}[1]{\foreignlanguage{hebrew}{#1}}

% --- Figure --------------------------------------------------------------------------
% \figw{larghezza in frazione di \textwidth}{file}{didascalia}
\newcommand{\figw}[3]{\begin{figure}[H]\centering\includegraphics[width=#1\textwidth,height=.70\textheight,keepaspectratio]{#2}\caption{#3}\end{figure}}
% \figh{altezza}{file1,file2,...}{didascalia}: immagini affiancate alla stessa altezza
\newcommand{\figitem}[2]{\includegraphics[height=#1,width=.31\textwidth,keepaspectratio]{#2}\hspace{.6em}}
\newcommand{\figh}[3]{\begin{figure}[H]\centering\forcsvlist{\figitem{#1}}{#2}\unskip\caption{#3}\end{figure}}

% --- Indice ----------------------------------------------------------------------------
\newcommand{\aut}[1]{\index[indice]{#1}}
\newcommand{\idx}[1]{\index[indice]{#1}}

\title{\Huge\bfseries @@CORSO@@\\[0.6em]\Large\mdseries Appunti del corso di @@DOCENTE@@\\ @@UNIVERSITA@@}
\author{}
\date{}

\begin{document}
\frontmatter
\maketitle

\chapter*{Come leggere questi appunti}
Il libro organizza gli argomenti secondo i prerequisiti e i collegamenti del corso. Le fonti e gli accorpamenti sono tracciati nell'inventario; l'ordine delle lezioni viene mantenuto se richiesto. Questi segnali distinguono la provenienza di ciò che leggi.

\begin{correzione}
Errore o incongruenza nel materiale. Il testo originale resta riportato, accanto alla correzione e al motivo.
\end{correzione}
\begin{dubbio}
Punto che le fonti del corso non permettono di risolvere: controllalo prima di usarlo all'esame.
\end{dubbio}
\begin{dallibro}
Spiegazione tratta dal libro di testo, con il numero di pagina. Chiarisce le slide senza aggiungere argomenti.
\end{dallibro}
\soloappunti
\noindent Così è segnalata una sezione senza slide, ricostruita soltanto dagli appunti in classe. \integrazione{così è segnata un'aggiunta di conoscenza generale, che non viene dal materiale del corso.}

\medskip
\noindent In fondo trovi il glossario dei concetti e l'indice analitico.

\tableofcontents

\mainmatter
% Campionario tecnico ed editoriale: gli ambienti qui mostrati sono facoltativi.
% Pagina di prova del preambolo: aggiungi un campione delle scritture e dei simboli
% del corso, compila, guarda il risultato, poi cancella il file e la sua riga \input.
\chapter{Prova tecnica}\label{cap:prova}

\section{Scritture}
Italiano con accenti: è, à, ì, ò, ù, perché, città. Parola inglese: \en{translation}\aut{Autore, Nome}\idx{concetto}.
Greco: \gr{παρθένος}, \gr{Ἰδοὺ ἡ παρθένος ἐν γαστρὶ ἕξει}. Rimando: capitolo~\ref{cap:prova}.

\begin{cit}[english]
A quotation in the language of the source.
\fontecit{(Fonte, anno)}
\end{cit}

\afronte[english]{The source text stands on the left, in its language, across several lines of the column.}{La traduzione sta a destra, su più righe della colonna.}

\section{Riquadri}
\begin{correzione}Prova di correzione.\end{correzione}
\begin{dubbio}Prova di dubbio.\end{dubbio}
\begin{dallibro}Prova dal libro (p.~1).\end{dallibro}
\soloappunti
Testo. \integrazione{prova di integrazione.}

\section{Tabelle, formule, codice, diagrammi}
\begin{center}\small
\begin{tabularx}{\linewidth}{@{}lYY@{}}
\toprule
& \textbf{Prima colonna} & \textbf{Seconda colonna} \\
\midrule
Riga & un testo abbastanza lungo da andare a capo dentro la cella della tabella & un altro testo abbastanza lungo da andare a capo \\
\bottomrule
\end{tabularx}
\end{center}

Formula in linea $E = mc^2$, $\alpha \to \beta$, e in evidenza:
\begin{align}
\sum_{i=1}^{n} i &= \frac{n(n+1)}{2} \\
\int_0^1 x^2 \, dx &= \frac{1}{3}
\end{align}

\begin{lstlisting}[numbers=left]
int main(void) { /* commento con è */
    return 0;
}
\end{lstlisting}

\begin{center}
\begin{tikzpicture}[node distance=2.4cm,every node/.style={draw,rounded corners,font=\small,inner sep=5pt}]
\node (a) {Sorgente};
\node[right=of a] (b) {Traduzione};
\draw[-{Latex}] (a) -- (b);
\end{tikzpicture}
\end{center}

\begin{insintesi}
Prova del riquadro di fine capitolo.
\end{insintesi}

% \input{capitoli/01-titolo}

\backmatter
\chapter{Glossario dei concetti}

Le voci rimandano al capitolo in cui il concetto è spiegato.

\begin{description}\small
\item[Concetto] Definizione breve (cap.~\ref{cap:prova}).
\end{description}

\clearpage\gdef\ultimocapitolo{Indice analitico}\gdef\ultimasezione{}
\printindex[indice]
\end{document}
